\begin{table}[H]
\scriptsize
\setlength{\tabcolsep}{2.5pt}
\renewcommand{\arraystretch}{0.84}

\caption{\label{tab:country-heterogeneity-counterfactual-gap}Household price policies: fiscal cost and effects on average inflation, individual P20--P80 inflation dispersion, and Q1--Q5 inflation inequality, 2021--2023}
\centering
\resizebox{\textwidth}{!}{%
\begin{tabular}{l*{10}{r}}
\toprule
Country & \shortstack{Fiscal cost\\(\% 2022 GDP)} & \multicolumn{9}{c}{Effect of price policies on}\\
\cmidrule(lr){3-11}
 & & \multicolumn{3}{c}{Average inflation} & \multicolumn{3}{c}{Individual P20--P80} & \multicolumn{3}{c}{Income Q1--Q5}\\
\cmidrule(lr){3-5} \cmidrule(lr){6-8} \cmidrule(lr){9-11}
 & & 2021 & 2022 & 2023 & 2021 & 2022 & 2023 & 2021 & 2022 & 2023\\
\midrule
Germany & 1.2$^{*}$ & 0.0 & -0.7 & -1.5 & 0.1 & -0.5 & -1.1 & 0.0 & -0.3 & -1.2\\
France & 2.9 & -0.1 & -2.7 & -0.6 & -0.1 & -2.7 & -1.1 & 0.0 & -1.1 & -0.5\\
Italy & 1.6 & -0.2 & -2.2 & 0.8 & -- & -- & -- & -0.1 & -0.5 & 0.0\\
Spain & 2.0$^{*}$ & -0.5 & -2.3 & -0.5 & -0.5 & -2.2 & -0.2 & -0.3 & -0.9 & -0.4\\
Netherlands & 1.0$^{*}$ & 0.0 & -2.5 & -0.5 & 0.0 & -2.5 & 0.3 & 0.0 & -1.2 & -0.3\\
\addlinespace
Belgium & 0.7$^{*}$ & 0.0 & -2.4 & 0.2 & 0.0 & -3.1 & -0.2 & 0.0 & -1.1 & 0.0\\
Austria & 0.9$^{*}$ & 0.0 & -0.2 & -0.8 & 0.0 & -0.2 & -0.7 & 0.0 & -0.1 & -0.4\\
Portugal & 1.9 & 0.0 & -1.0 & -2.5 & -- & -- & -- & 0.0 & -0.5 & -1.7\\
Greece & 2.4$^{*}$ & -0.3 & -0.9 & 0.3 & -0.3 & -0.6 & -0.3 & -0.2 & -0.3 & 0.0\\
Finland & 0.2 & 0.0 & -0.4 & -0.4 & 0.0 & -0.5 & 0.0 & 0.0 & 0.2 & 0.1\\
\addlinespace
Ireland & 0.5$^{*}$ & 0.0 & -0.9 & 0.0 & 0.0 & -1.1 & -0.3 & 0.0 & -0.5 & -0.1\\
Slovakia & 1.7 & 0.0 & 0.0 & -3.6 & 0.0 & 0.0 & -2.7 & 0.0 & 0.0 & -2.0\\
Croatia & 0.9 & 0.0 & -1.2 & 0.4 & 0.0 & -0.7 & -0.4 & 0.0 & -1.1 & 0.5\\
Lithuania & 0.8 & 0.0 & -0.4 & -0.1 & 0.0 & -0.5 & 0.0 & 0.0 & -0.4 & -0.1\\
Slovenia & 0.3 & 0.0 & -1.0 & -0.1 & 0.0 & -0.6 & -0.8 & 0.0 & -0.1 & -0.6\\
\addlinespace
Luxembourg & 0.8$^{*}$ & 0.0 & -0.7 & -0.6 & 0.0 & -0.6 & -0.4 & 0.0 & -0.3 & -0.2\\
Latvia & 1.1 & 0.0 & -1.5 & -0.7 & 0.0 & -1.8 & -0.3 & 0.0 & -1.7 & -0.5\\
Estonia & 0.5$^{*}$ & -0.2 & -1.1 & 0.4 & -0.2 & -1.8 & -0.1 & -0.1 & -0.9 & 0.4\\
Cyprus & 0.7 & -0.1 & -0.6 & 0.6 & 0.0 & -0.4 & -0.1 & 0.0 & -0.2 & 0.2\\
Malta & 2.6$^{*}$ & 0.0 & -0.9 & 0.0 & 0.0 & -0.1 & 0.1 & 0.0 & -0.1 & 0.0\\
\addlinespace
\textbf{Euro area} & \textbf{1.6} & \textbf{-0.1} & \textbf{-1.7} & \textbf{-0.6} & \textbf{-0.1} & \textbf{-1.6} & \textbf{-0.7} & \textbf{-0.1} & \textbf{-0.6} & \textbf{-0.6}\\
\bottomrule
\end{tabular}%
}
\begin{flushleft}
\scriptsize{Notes. Effects are observed minus counterfactual annual inflation, in percentage points. The first block reports average inflation. The second reports the P80 minus P20 range of household inflation; national endpoints use HBS survey weights. Italy and Portugal are unavailable at the required household-product detail and are shown as dashes; the euro-area row averages national endpoints over the other 18 countries using annual HICP country weights. The third block is the Q1 minus Q5 gap for all 20 countries. Costs are scheme budgets as a percentage of 2022 GDP, counting shared envelopes once. An asterisk marks costs including business beneficiaries; the household share is not identified. A dash denotes incomplete coverage, not zero. France's fuel rebate excludes EUR 1.6 billion paid to road freight; its electricity shield also covers microbusinesses. Source periods vary (France and Lithuania: 2022); Slovak amounts are gross 2023 expenditure. Dutch quintile effects measure basket exposure, not cap eligibility. Sources: Eurostat, Bruegel, national sources, authors' calculations.}
\end{flushleft}
\end{table}
